\documentclass[12pt,a4paper]{article}

% === ENCODAGE ET LANGUE ===
\usepackage[utf8]{inputenc}
\usepackage[T1]{fontenc}
\usepackage[french]{babel}

% === MISE EN PAGE ===
\usepackage{geometry}
\geometry{margin=2.5cm, headheight=14.5pt}
\usepackage{setspace}
\onehalfspacing

% === POLICES ET TYPOGRAPHIE ===
\usepackage{lmodern}
\usepackage{microtype}

% === GRAPHIQUES ET FIGURES ===
\usepackage{graphicx}
\graphicspath{{figures/}}
\usepackage{tikz}
\usetikzlibrary{arrows.meta, shapes.geometric, positioning, calc, fit, backgrounds}
\usepackage{pgfplots}
\pgfplotsset{compat=1.18}
\usepackage{float}

% === TABLEAUX ===
\usepackage{booktabs}
\usepackage{multirow}
\usepackage{tabularx}
\usepackage{array}
\usepackage{colortbl}
\usepackage{makecell}

% === COULEURS ===
\usepackage[table]{xcolor}
\definecolor{upsred}{RGB}{180, 30, 40}
\definecolor{criticalred}{RGB}{220, 53, 69}
\definecolor{importantyellow}{RGB}{255, 193, 7}
\definecolor{moderategray}{RGB}{108, 117, 125}
\definecolor{lightgray}{RGB}{245, 245, 245}
\definecolor{codebg}{RGB}{248, 249, 250}

% === MATHS ===
\usepackage{amsmath}
\usepackage{amssymb}

% === LISTINGS (CODE) ===
\usepackage{listings}
\lstset{
  basicstyle=\ttfamily\small,
  backgroundcolor=\color{codebg},
  frame=single,
  rulecolor=\color{gray!30},
  breaklines=true,
  columns=fullflexible,
  literate={é}{{\'e}}1 {è}{{\`e}}1 {ê}{{\^e}}1 {à}{{\`a}}1 {ù}{{\`u}}1 {û}{{\^u}}1 {ç}{{\c{c}}}1 {ô}{{\^o}}1 {î}{{\^i}}1
}

% === LIENS ===
\usepackage{hyperref}
\hypersetup{
  colorlinks=true,
  linkcolor=upsred,
  citecolor=upsred,
  urlcolor=blue!70!black,
  pdfauthor={Mohamed DIALLO et al.},
  pdftitle={GeoR2LLMs - Rapport Chef d'oeuvre},
  pdfsubject={Systeme Question-Reponse Geographique avec RAG}
}

% === BIBLIOGRAPHIE ===
\usepackage{cite}

% === EN-TETES ET PIEDS DE PAGE ===
\usepackage{fancyhdr}
\pagestyle{fancy}
\fancyhf{}
\fancyhead[L]{\small\textit{GeoR2LLMs - Rapport Chef d'\oe uvre}}
\fancyhead[R]{\small\textit{M2 IA - UPS}}
\fancyfoot[C]{\thepage}
\renewcommand{\headrulewidth}{0.4pt}
\renewcommand{\footrulewidth}{0.4pt}

\fancypagestyle{plain}{
  \fancyhf{}
  \fancyfoot[C]{\thepage}
  \renewcommand{\headrulewidth}{0pt}
  \renewcommand{\footrulewidth}{0pt}
}

% === TITRES ===
\usepackage{titlesec}
\titleformat{\section}{\Large\bfseries\color{upsred}}{\thesection}{1em}{}
\titleformat{\subsection}{\large\bfseries}{\thesubsection}{1em}{}
\titleformat{\subsubsection}{\normalsize\bfseries\itshape}{\thesubsubsection}{1em}{}

% === ENUMERATION ===
\usepackage{enumitem}

% === CAPTION ===
\usepackage{caption}
\usepackage{subcaption}
\captionsetup{font=small, labelfont=bf, format=hang}

% ====================================================================
\begin{document}

% ====================================================================
% PAGE DE GARDE
% ====================================================================
\begin{titlepage}
  \centering
  \vspace*{0.5cm}

  
\includegraphics[width=5cm]{Ressources/logo-UT-site.png}\\[0.8cm]

  {\Large\textsc{Universit\'e Toulouse III - Paul Sabatier}}\\[0.3cm]
  {\large\textsc{Master 2 Intelligence Artificielle}}\\[0.2cm]
  {\normalsize IRIT - Institut de Recherche en Informatique de Toulouse}\\[1.5cm]

  \rule{\textwidth}{1.5pt}\\[0.6cm]
  {\huge\bfseries Syst\`eme Question-R\'eponse\\[0.3cm] en G\'eographie avec RAG}\\[0.4cm]
  {\Large\textit{Rapport Chef d'\oe uvre}}\\[0.3cm]
  {\large Projet GeoR2LLMs}\\[0.4cm]
  \rule{\textwidth}{1.5pt}\\[2cm]

  \begin{minipage}{0.45\textwidth}
    \begin{flushleft}
      \textbf{R\'ealis\'e par}\\[0.3cm]
      Mohamed DIALLO\\
      Hiba L'MOUDDEN\\
      Feras MAKHLOUF\\
      Osama BAAZZI\\
      Ousama BOUTEZROUT
    \end{flushleft}
  \end{minipage}
  \hfill
  \begin{minipage}{0.45\textwidth}
    \begin{flushright}
      \textbf{Sous la direction de}\\[0.3cm]
      Madame Lynda TAMINE-LECHANI\\
      Professeure, IRIT\\[0.5cm]
      \textbf{UE} Chef d'\oe uvre\\
      \textbf{Ann\'ee} 2025-2026
    \end{flushright}
  \end{minipage}

  \vfill
  {\large Mars 2026}
\end{titlepage}

% ====================================================================
% TABLE DES MATIERES
% ====================================================================
\newpage
\tableofcontents
\newpage

% ====================================================================
% 1. INTRODUCTION
% ====================================================================
\section{Introduction}
\label{sec:introduction}

\subsection{Contexte g\'en\'eral}

Ce rapport constitue le livrable final du projet Chef d'\oe uvre intitul\'e \og Syst\`eme Question-R\'eponse en G\'eographie \fg, r\'ealis\'e dans le cadre du Master~2 Intelligence Artificielle de l'Universit\'e Toulouse III - Paul Sabatier, sous la direction de Madame Lynda Tamine-Lechani \`a l'IRIT. Le projet s'inscrit dans le cadre du projet europ\'een GeoR2LLMs (\textit{Knowledgeable and Multimodal Geographic Large Language Models Grounded with Reasoning and Retrieval}).

L'\'etude bibliographique r\'ealis\'ee au premier semestre (septembre-d\'ecembre 2025) a permis d'identifier les d\'efis sp\'ecifiques du QA g\'eographique~\cite{mai_geographic_2021}, les limites des LLMs pour le raisonnement spatial~\cite{cohn_evaluating_2024}, et le potentiel de l'architecture RAG~\cite{lewis2020retrieval, gao2023ragsurvey} pour compenser les lacunes de connaissance des mod\`eles. Le pr\'esent rapport couvre la p\'eriode janvier-mars 2026, consacr\'ee \`a l'\textbf{impl\'ementation et l'\'evaluation} d'un syst\`eme QA en configuration RAG.

\subsection{Objectifs}

L'objectif de cette phase est d'impl\'ementer et d'\'evaluer les performances d'un syst\`eme d'acc\`es \`a l'information (question-r\'eponse) dans une configuration de g\'en\'eration (G) augment\'ee par la recherche d'information (R). Les exigences fix\'ees par le cahier des charges sont les suivantes\,:

\begin{itemize}[leftmargin=*]
  \item La fonction R doit \^etre fond\'ee sur des mod\`eles de recherche d'information de l'\'etat de l'art, avec des biblioth\`eques reconnues.
  \item La fonction G doit reposer sur des mod\`eles de diff\'erentes architectures et de diff\'erentes tailles.
  \item Le syst\`eme doit pouvoir \^etre d\'eploy\'e avec diff\'erents corpus et diff\'erentes ressources externes.
  \item Les m\'ethodes doivent \^etre r\'ef\'erenc\'ees dans des articles scientifiques.
  \item L'\'evaluation exp\'erimentale doit se fonder sur des standards reconnus, \`a savoir la d\'efinition de baselines sans RAG, la d\'efinition de configurations RAG, l'utilisation de corpus standards et l'application de protocoles d'\'evaluation document\'es.
\end{itemize}

\subsection{Organisation du travail}

Le travail a \'et\'e structur\'e en deux \'etapes compl\'ementaires\,:

\begin{itemize}[leftmargin=*]
  \item \textbf{\'Etape 1 - Plateforme de base.} Mise en place d'un pipeline RAG avec un seul LLM (Qwen3-0.6B), un seul encodeur (BGE-m3), un seul dataset (GeoSQA) et Wikipedia fran\c{c}ais comme corpus. L'objectif est d'\'etablir un pipeline fonctionnel et une baseline de r\'ef\'erence.
  \item \textbf{\'Etape 2 - Exploration.} Extension de l'\'evaluation \`a 5 LLMs (dont 4 exploitables), 3 datasets g\'eographiques (GeoSQA, GKMC, Tourism-QA) et une analyse factorielle crois\'ee des performances.
\end{itemize}

Le code source et les notebooks sont disponibles sur le d\'ep\^ot GitHub du projet\footnote{\url{https://github.com/Hiba-Lmoudden/GeoR2LLMs}}.

\subsection{Organisation du rapport}

La section~\ref{sec:rappel_biblio} rappelle bri\`evement les fondements th\'eoriques issus de l'\'etude bibliographique. La section~\ref{sec:architecture} d\'ecrit l'architecture du syst\`eme. La section~\ref{sec:methodologie} d\'etaille la m\'ethodologie exp\'erimentale. Les sections~\ref{sec:phase1} et~\ref{sec:phase2} rapportent les r\'esultats des deux \'etapes. La section~\ref{sec:analyse} pr\'esente l'analyse factorielle crois\'ee. La section~\ref{sec:discussion} propose une discussion critique. Enfin, la section~\ref{sec:conclusion} conclut et ouvre des perspectives.


% ====================================================================
% 2. RAPPEL DES FONDEMENTS THEORIQUES
% ====================================================================
\section{Rappel des fondements th\'eoriques}
\label{sec:rappel_biblio}

Cette section synth\'etise les \'el\'ements cl\'es de l'\'etude bibliographique r\'ealis\'ee au premier semestre, n\'ecessaires \`a la compr\'ehension des choix d'impl\'ementation et des r\'esultats exp\'erimentaux.

\subsection{Grands mod\`eles de langue (LLMs)}

Les LLMs sont des mod\`eles de langue pr\'e-entra\^in\'es fond\'es sur l'architecture Transformer~\cite{vaswani2017attention}, contenant des dizaines \`a des centaines de milliards de param\`etres~\cite{zhao_survey_2025}. L'architecture Transformer repose sur le m\'ecanisme d'\textbf{attention multi-t\^etes} (\textit{multi-head self-attention}), qui permet \`a chaque token de la s\'equence de prendre en compte l'ensemble des tokens pr\'ec\'edents via des fonctions de score pond\'er\'ees. Formellement, pour une t\^ete d'attention $h$, la sortie est $\text{Attention}(Q, K, V) = \text{softmax}\!\left(\frac{QK^\top}{\sqrt{d_k}}\right)V$, o\`u $Q$, $K$, $V$ sont les projections lin\'eaires de l'entr\'ee. Ce m\'ecanisme remplace les r\'ecurrences (RNN/LSTM) et les convolutions, permettant une parall\'elisation compl\`ete de l'entra\^inement.

On distingue trois familles architecturales. Les \textbf{encodeurs} (BERT~\cite{devlin2019bert}) utilisent une attention bidirectionnelle et sont adapt\'es \`a la classification et \`a l'extraction. Les \textbf{d\'ecodeurs} (GPT, LLaMA~\cite{touvron2023llama}, Qwen~\cite{yang2025qwen3}) utilisent une attention causale (masqu\'ee) et g\'en\`erent du texte de mani\`ere autor\'egressive, token par token. Les mod\`eles \textbf{encodeur-d\'ecodeur} (T5) combinent les deux m\'ecanismes. Notre projet utilise exclusivement des d\'ecodeurs, car la t\^ache de g\'en\'eration de r\'eponses (QA g\'en\'eratif) requiert une sortie textuelle libre.

Au-del\`a d'un certain seuil de taille, ces mod\`eles exhibent des capacit\'es \'emergentes~\cite{wei2022emergent} telles que l'apprentissage en contexte (\textit{in-context learning}), la capacit\'e de raisonnement par \'etapes (\textit{Chain-of-Thought}~\cite{wei_chain--thought_2022}), et l'instruction following. L'alignement par RLHF (\textit{Reinforcement Learning from Human Feedback})~\cite{ouyang2022training} permet ensuite d'affiner le mod\`ele pour qu'il suive les instructions de mani\`ere fiable et produise des r\'eponses utiles, ce qui est fondamental pour les mod\`eles de type Instruct utilis\'es dans ce projet.

Dans le cadre de ce projet, trois familles de LLMs ont \'et\'e consid\'er\'ees. La famille \textbf{Qwen3} (Alibaba)~\cite{yang2025qwen3} avec les tailles 0.6B et 1.7B param\`etres. La famille \textbf{Llama 3.2} (Meta)~\cite{meta2024llama32} en version Instruct avec les tailles 1B et 3B. Enfin, la famille \textbf{Mistral} avec le mod\`ele Ministral-3-3B-Instruct-2512, qui n'a finalement pas pu \^etre charg\'e en raison d'une incompatibilit\'e technique (cf. section~\ref{sec:ministral}). Tous ces mod\`eles sont de taille inf\'erieure \`a 3B param\`etres, une contrainte impos\'ee par le GPU T4 (16~Go VRAM) disponible sur Google Colab.

\subsection{G\'en\'eration augment\'ee par la recherche d'information (RAG)}

L'architecture RAG, introduite par Lewis et al.~\cite{lewis2020retrieval}, propose d'augmenter les LLMs avec des connaissances externes r\'ecup\'er\'ees dynamiquement depuis un corpus. L'intuition fondamentale est que les LLMs, m\^eme de grande taille, ne peuvent pas stocker l'ensemble des connaissances factuelles dans leurs param\`etres~\cite{petroni2019language}. Le RAG contourne cette limitation en externalisant une partie de la m\'emoire factuelle dans un corpus interrogeable.

\'Etant donn\'ee une question $q$, un module de recherche (retriever) identifie les $K$ passages les plus pertinents dans un corpus $\mathcal{C}$. La recherche repose sur le \textbf{Dense Passage Retrieval}~\cite{karpukhin2020dense}, o\`u la question et les passages sont encod\'es dans un m\^eme espace vectoriel dense par un mod\`ele d'embeddings. La similarit\'e est alors calcul\'ee par produit scalaire ou cosinus entre les vecteurs. Les $K$ passages les plus similaires sont ensuite concat\'en\'es avec la question dans le prompt du g\'en\'erateur, qui produit la r\'eponse en se fondant \`a la fois sur ses connaissances internes et sur le contexte fourni.

Gao et al.~\cite{gao2023ragsurvey} distinguent le \textbf{RAG na\"if} (\textit{retrieve-then-read}), qui injecte directement les documents r\'ecup\'er\'es dans le prompt, du \textbf{RAG avanc\'e}, qui ajoute des \'etapes de r\'e\'ecriture de requ\^ete, de re-ranking par cross-encoder, ou de RAG it\'eratif. Nous avons impl\'ement\'e un RAG na\"if, la configuration la plus directe et la plus reproductible, ce qui constitue \`a la fois une force (reproductibilit\'e, simplicit\'e d'analyse) et une limitation (absence d'optimisation du retrieval). Liu et al.~\cite{liu2023lost} ont par ailleurs montr\'e que les LLMs peinent \`a exploiter les informations situ\'ees au milieu de longs contextes (\textit{lost in the middle}), un ph\'enom\`ene pertinent pour le choix du nombre de documents r\'ecup\'er\'es ($K = 5$ dans notre cas).

\subsection{Question-R\'eponse en g\'eographie}

Le QA g\'eographique se distingue du QA g\'en\'eraliste par plusieurs sp\'ecificit\'es identifi\'ees par Mai et al.~\cite{mai_geographic_2021}. Les questions g\'eographiques requi\`erent un raisonnement spatial (proximit\'e, inclusion, adjacence), la compr\'ehension de relations topologiques (``au nord de'', ``travers\'e par''), et l'int\'egration de connaissances sp\'ecifiques au domaine (capitales, altitudes, populations). Mai et al. proposent une taxonomie distinguant les questions factuelles (``Quelle est la capitale du Burkina Faso~?''), les questions de raisonnement spatial (``Quel pays se situe entre la France et l'Espagne~?'') et les questions de comparaison (``Quel fleuve est le plus long~?'').

Cohn et Blackwell~\cite{cohn_evaluating_2024} d\'emontrent qu'aucun LLM actuel ne parvient \`a raisonner de mani\`ere fiable sur les directions cardinales, m\^eme avec une temp\'erature de z\'ero. Ramrakhiyani et al.~\cite{ramrakhiyani_zero-shot_2023} montrent que les mod\`eles de petite taille peinent \`a restituer des faits g\'eographiques pr\'ecis en \'evaluation \textit{zero-shot}, avec des performances variant fortement selon le type de connaissance (les capitales sont mieux retenues que les coordonn\'ees). Ces constats justifient l'int\'er\^et d'\'etudier l'apport du RAG pour compenser ces lacunes, en fournissant aux mod\`eles des passages contenant les faits g\'eographiques n\'ecessaires plut\^ot que de compter exclusivement sur leur m\'emoire param\'etrique.


% ====================================================================
% 3. ARCHITECTURE DU SYSTEME
% ====================================================================
\section{Architecture du syst\`eme}
\label{sec:architecture}

Le syst\`eme GeoR2LLMs repose sur une architecture RAG na\"ive (\textit{retrieve-then-read}) compos\'ee de trois modules. Un module \textbf{Retriever} qui r\'ecup\`ere les passages pertinents, un module \textbf{Generator} qui produit la r\'eponse, et un module \textbf{Evaluator} qui mesure les performances. La figure~\ref{fig:architecture} illustre le pipeline complet.

\begin{figure}[H]
\centering
\resizebox{\textwidth}{!}{%
\begin{tikzpicture}[
  node distance=1.2cm and 2cm,
  box/.style={rectangle, draw, rounded corners=4pt, minimum height=1.2cm, 
              minimum width=3cm, align=center, font=\small},
  data/.style={box, fill=blue!8, draw=blue!40},
  process/.style={box, fill=green!8, draw=green!40},
  model/.style={box, fill=orange!10, draw=orange!40},
  result/.style={box, fill=red!8, draw=red!40},
  arrow/.style={-{Stealth[length=3mm]}, thick, draw=gray!70},
  lbl/.style={font=\scriptsize\itshape, text=gray}
]

% Donnees d'entree
\node[data] (wiki) {Corpus\\{\scriptsize Wikipedia / Reviews}};
\node[data, below=1.5cm of wiki] (question) {Question\\{\scriptsize GeoSQA / GKMC / Tourism-QA}};

% Retriever (ligne du haut)
\node[process, right=2cm of wiki] (chunking) {Chunking\\{\scriptsize 512 car., overlap 50}};
\node[model, right=1.5cm of chunking] (bge) {BGE-m3\\{\scriptsize 1024 dim.}};
\node[process, right=1.5cm of bge] (faiss) {FAISS\\{\scriptsize IndexFlatIP}};

% Ligne du bas : Encodage aligné sous chunking, Top-K aligné sous faiss
\node[model] (embed_q) at (chunking |- question) {Encodage\\{\scriptsize question}};
\node[process] (topk) at (faiss |- question) {Top-K\\{\scriptsize K=5}};

% LLM aligné sous Top-K
\node[model, below=1.5cm of topk] (llm) {LLM\\{\scriptsize Generator}};
\node[result, right=1.5cm of llm] (reponse) {R\'eponse\\{\scriptsize g\'en\'er\'ee}};

% Fleches corpus (ligne haute)
\draw[arrow] (wiki) -- (chunking);
\draw[arrow] (chunking) -- (bge);
\draw[arrow] (bge) -- node[above, lbl] {embeddings} (faiss);

% Fleches question (ligne basse)
\draw[arrow] (question) -- (embed_q);
\draw[arrow] (embed_q) -- node[above, lbl] {vecteur requ\^ete} (topk);
\draw[arrow] (faiss) -- node[right, lbl] {index} (topk);

% Fleches generation
\draw[arrow] (topk) -- node[right, lbl] {contexte} (llm);
\draw[arrow] (question.south) |- (llm.west);
\draw[arrow] (llm) -- (reponse);

% Cadres
\begin{scope}[on background layer]
  \node[fit=(chunking)(bge)(faiss), 
        draw=green!60, dashed, rounded corners=8pt, inner sep=10pt, 
        label={[font=\small\bfseries, text=green!60!black]above:Module Retriever}] {};
  \node[fit=(llm)(reponse), 
        draw=orange!60, dashed, rounded corners=8pt, inner sep=10pt, 
        label={[font=\small\bfseries, text=orange!60!black]below:Module Generator}] {};
\end{scope}

\end{tikzpicture}
}
\caption{Architecture du pipeline RAG GeoR2LLMs. Le corpus est d\'ecoup\'e en chunks, 
encod\'e par BGE-m3 et index\'e dans FAISS. Les questions sont encod\'ees, puis les 
top-K passages r\'ecup\'er\'es sont fournis comme contexte au g\'en\'erateur LLM.}
\label{fig:architecture}
\end{figure}

\subsection{Module Retriever}

Le retriever utilise le mod\`ele d'embeddings multilingue \textbf{BGE-m3} (BAAI), pr\'esent\'e par Chen et al.~\cite{chen2024bgem3} aux Findings d'ACL 2024. Ce mod\`ele produit des vecteurs denses de 1024 dimensions et supporte plus de 100 langues, ce qui le rend adapt\'e \`a notre pipeline o\`u les questions sont en anglais et le corpus peut \^etre en fran\c{c}ais ou en anglais. Le choix de BGE-m3 r\'epond \`a l'exigence du cahier des charges d'utiliser un \textit{sentence transformer} reconnu.

Les documents du corpus sont d\'ecoup\'es en \textit{chunks} de 512 caract\`eres avec un chevauchement de 50 caract\`eres pour pr\'eserver le contexte aux fronti\`eres. Les embeddings des chunks sont stock\'es dans un index \textbf{FAISS}~\cite{johnson2019faiss} de type \texttt{IndexFlatIP}, qui effectue une recherche par produit scalaire apr\`es normalisation L2 des vecteurs. Cette op\'eration est math\'ematiquement \'equivalente \`a la similarit\'e cosinus. Pour une requ\^ete $q$, le retriever retourne les $K$ passages les plus similaires.

\begin{equation}
  \text{Top-K}(q) = \underset{d_i \in \mathcal{C}}{\text{argmax-K}} \; \cos(\mathbf{e}_q, \mathbf{e}_{d_i})
\end{equation}

\subsection{Module Generator}

En \'etape 1, le g\'en\'erateur est \textbf{Qwen3-0.6B}~\cite{yang2025qwen3}, un mod\`ele de 0.6 milliard de param\`etres. En \'etape 2, le syst\`eme est \'etendu \`a 4 mod\`eles gr\^ace \`a un composant \texttt{UnifiedGenerator} qui encapsule le chargement, le formatage de prompt et la g\'en\'eration pour les trois familles (Qwen, Llama, Mistral). La lib\'eration m\'emoire GPU explicite entre chaque mod\`ele (\texttt{del model}, \texttt{torch.cuda.empty\_cache()}) permet d'\'evaluer s\'equentiellement plusieurs LLMs sur un GPU T4 de 16~Go sans erreur de d\'epassement de m\'emoire.

Le prompt suit la structure suivante.

\begin{lstlisting}[language=Python, caption={Prompt template pour le mode RAG}]
system: "Tu es un expert en geographie. Reponds
         uniquement avec la lettre correcte."
user:   "Contexte: {documents recuperes}
         Question: {question}
         Options: A) ... B) ... C) ... D) ..."
\end{lstlisting}

En mode \textbf{Baseline}, aucun contexte n'est fourni et le mod\`ele r\'epond uniquement \`a partir de ses connaissances internes. Cette configuration sert de point de comparaison pour mesurer l'apport du RAG.

\subsection{Module Evaluator}

L'\'evaluateur mesure la performance selon deux axes. D'une part, la qualit\'e de g\'en\'eration (Accuracy MCQ, Exact Match, F1 token-level), conform\'ement aux m\'etriques standard du QA~\cite{rajpurkar2016squad}. D'autre part, la qualit\'e du retrieval (Hit Rate@K, MRR@K), selon les conventions du Dense Passage Retrieval~\cite{karpukhin2020dense}. Pour Tourism-QA, un \texttt{POIEvaluator} d\'edi\'e calcule l'Accuracy@1, le Hit Rate@3 et le F1 sur le nom du POI.


% ====================================================================
% 4. METHODOLOGIE
% ====================================================================
\section{M\'ethodologie}
\label{sec:methodologie}

\subsection{Mod\`eles \'evalu\'es}

Le tableau~\ref{tab:modeles} r\'ecapitule les 5 mod\`eles pr\'evus par le cahier des charges. Seuls 4 d'entre eux ont produit des r\'esultats exploitables (cf. section~\ref{sec:ministral}).

\begin{table}[H]
\centering
\caption{Mod\`eles LLM \'evalu\'es. Les param\`etres de g\'en\'eration suivent les recommandations de chaque \'equipe de d\'eveloppement.}
\label{tab:modeles}
\small
\begin{tabular}{@{} l l c c l @{}}
\toprule
\textbf{Mod\`ele} & \textbf{Famille} & \textbf{Taille} & \textbf{Temp.}
  & \textbf{Identifiant HuggingFace} \\
\midrule
Qwen3-0.6B   & Qwen    & 0.6B & 0.7 & \texttt{Qwen/Qwen3-0.6B} \\
\rowcolor{lightgray}
Qwen3-1.7B   & Qwen    & 1.7B & 0.7 & \texttt{Qwen/Qwen3-1.7B} \\
Llama-3.2-1B & Llama   & 1.0B & 0.6 & \texttt{meta-llama/Llama-3.2-1B-Instruct} \\
\rowcolor{lightgray}
Llama-3.2-3B & Llama   & 3.0B & 0.6 & \texttt{meta-llama/Llama-3.2-3B-Instruct} \\
Ministral-3B & Mistral & 3.0B & 0.7 & \texttt{mistralai/Ministral-3-3B-Instruct-2512} \\
\bottomrule
\end{tabular}
\end{table}

Les mod\`eles de plus grande taille mentionn\'es comme optionnels dans le cahier des charges (Qwen3-8B, Llama-3.1-8B-Instruct, Ministral-3-8B-Instruct-2512) n'ont pas pu \^etre \'evalu\'es car leur empreinte m\'emoire en float16 d\'epasse les 15.6~Go disponibles sur le GPU T4 une fois BGE-m3 (2.27~Go) charg\'e.

\subsection{Datasets et corpus}

\begin{table}[H]
\centering
\small
\caption{Datasets utilis\'es pour l'\'evaluation. Le corpus RAG diff\`ere selon le dataset.}
\label{tab:datasets}
\begin{tabular}{@{} l c c l l @{}}
\toprule
\textbf{Dataset} & \textbf{Questions} & \textbf{Langue} & \textbf{T\^ache} & \textbf{Corpus RAG} \\
\midrule
GeoSQA~\cite{huang_geosqa_2019}
  & 200 & ZH $\to$ EN & QCM (A/B/C/D) & Wikipedia FR \\
\rowcolor{lightgray}
GKMC & 200 & EN & QCM (A/B/C/D) & Wikipedia FR \\
Tourism-QA~\cite{tourismqa2021}
  & 100 & EN & Recommandation POI & Reviews utilisateurs \\
\bottomrule
\end{tabular}
\end{table}

\textbf{GeoSQA}~\cite{huang_geosqa_2019} est un benchmark de question-r\'eponse sc\'enaris\'e en g\'eographie, d\'evelopp\'e par Huang et al. et pr\'esent\'e \`a EMNLP-IJCNLP 2019. Il contient 1\,981 sc\'enarios et 4\,110 QCM au niveau lyc\'ee chinois. Les questions, originellement en chinois, ont \'et\'e traduites en anglais dans la version HuggingFace (\texttt{rfr2003/Geo\_Benchmark}). Le split de test contient 838 questions, dont 200 sont \'echantillonn\'ees al\'eatoirement pour nos \'evaluations.

\textbf{GKMC} (\textit{Geography Knowledge Multiple Choice}) est un benchmark de connaissances g\'eographiques g\'en\'erales, nativement en anglais, disponible dans le m\^eme d\'ep\^ot HuggingFace (\texttt{rfr2003/Geo\_Benchmark}, configuration \texttt{GKMC}). Contrairement \`a GeoSQA qui \'evalue le raisonnement spatial par sc\'enarios, GKMC teste des connaissances encyclop\'ediques factuelles (capitales, fleuves, reliefs, fronti\`eres). Les questions \'etant nativement en anglais, ce dataset permet de dissocier l'effet du mismatch linguistique (pr\'esent dans GeoSQA avec la traduction ZH$\to$EN) de la difficult\'e intrins\`eque de la t\^ache. Comme pour GeoSQA, 200 questions sont \'echantillonn\'ees al\'eatoirement depuis le split de test.

\textbf{Tourism-QA}~\cite{tourismqa2021} \'evalue la recommandation de points d'int\'er\^et (POI) \`a partir de requ\^etes touristiques. Contrairement aux deux premiers benchmarks, il ne s'agit pas d'un QCM mais d'une t\^ache ouverte dans laquelle le syst\`eme doit recommander le bon POI en se fondant sur les avis (reviews) d'utilisateurs, qui constituent \'egalement le corpus RAG.

Le \textbf{corpus Wikipedia} est constitu\'e de 300 articles charg\'es en streaming depuis le dump Wikipedia fran\c{c}ais (\texttt{wikimedia/wikipedia}, configuration \texttt{20231101.fr}). Apr\`es d\'ecoupage en chunks de 512 caract\`eres avec chevauchement de 50 caract\`eres, ce corpus produit environ 2\,527 chunks. Pour Tourism-QA, le corpus est constitu\'e des reviews d'utilisateurs associ\'ees aux POI du dataset, nativement en anglais.

\subsection{Configuration exp\'erimentale}

Toutes les exp\'eriences ont \'et\'e r\'ealis\'ees sur Google Colab avec un GPU NVIDIA Tesla T4 (16~Go VRAM). Les param\`etres communs sont les suivants. La graine al\'eatoire est fix\'ee \`a 42 pour la reproductibilit\'e. Le nombre de documents r\'ecup\'er\'es est $K = 5$. La taille des chunks est de 512 caract\`eres avec un chevauchement de 50 caract\`eres. Le nombre maximal de tokens g\'en\'er\'es est fix\'e \`a 100 (150 pour Tourism-QA). Chaque mod\`ele est \'evalu\'e dans deux conditions. En \textbf{Baseline}, le LLM r\'epond seul, sans contexte externe. En \textbf{RAG}, le LLM re\c{c}oit les top-5 documents r\'ecup\'er\'es comme contexte dans son prompt.

\subsection{M\'etriques d'\'evaluation}

\subsubsection{M\'etriques de g\'en\'eration}

Les m\'etriques de g\'en\'eration suivent les conventions de SQuAD~\cite{rajpurkar2016squad}.

\begin{itemize}[leftmargin=*]
  \item \textbf{Accuracy MCQ.} La premi\`ere lettre valide (A, B, C ou D) pr\'esente dans la r\'eponse g\'en\'er\'ee est extraite par expression r\'eguli\`ere, puis compar\'ee \`a la r\'ef\'erence. Cette m\'etrique est robuste \`a la verbosit\'e du mod\`ele.
  \begin{equation}
    \text{Accuracy} = \frac{1}{N} \sum_{i=1}^{N} \mathbb{1}[\hat{y}_i = y_i]
  \end{equation}

  \item \textbf{Exact Match (EM).} Correspondance exacte entre le texte pr\'edit normalis\'e et la r\'ef\'erence. Dans le contexte des QCM, cette m\'etrique est structurellement nulle car les mod\`eles g\'en\`erent des phrases compl\`etes alors que la r\'ef\'erence est une lettre seule.

  \item \textbf{F1 Score.} F1 au niveau des tokens entre la pr\'ediction et la r\'ef\'erence, apr\`es tokenisation et normalisation.
  \begin{equation}
    F_1 = \frac{2 \cdot P \cdot R}{P + R}, \quad P = \frac{|\hat{T} \cap T|}{|\hat{T}|}, \quad R = \frac{|\hat{T} \cap T|}{|T|}
  \end{equation}

  \item \textbf{Accuracy@1.} Pour Tourism-QA, proportion de cas o\`u le premier POI recommand\'e correspond au POI correct.
\end{itemize}

\subsubsection{M\'etriques de retrieval}

\begin{itemize}[leftmargin=*]
  \item \textbf{Hit Rate@K.} Proportion de questions pour lesquelles au moins un des $K$ documents r\'ecup\'er\'es contient le texte de la r\'eponse correcte.
  \begin{equation}
    \text{HR@K} = \frac{1}{N} \sum_{i=1}^{N} \mathbb{1}\left[\exists j \leq K : y_i \in d_j^{(i)}\right]
  \end{equation}

  \item \textbf{MRR@K} (\textit{Mean Reciprocal Rank}). Moyenne du rang r\'eciproque du premier document pertinent parmi les top-K.
  \begin{equation}
    \text{MRR@K} = \frac{1}{N} \sum_{i=1}^{N} \frac{1}{\text{rank}_i}
  \end{equation}
\end{itemize}

\subsubsection{Taux d'accroissement relatif}

Pour comparer Baseline et RAG, le taux d'accroissement relatif $\Delta_{rel}$ est calcul\'e pour chaque m\'etrique.

\begin{equation}
  \Delta_{rel} = \frac{v_{\text{RAG}} - v_{\text{Base}}}{v_{\text{Base}}} \times 100
\end{equation}

\subsection{Correspondance avec les métriques RAGAS}

Les métriques RAGAS n'ont pas pu être calculées directement pour trois raisons cumulatives : (1) RAGAS requiert un juge externe (GPT-4) dont l'API est payante et indisponible dans notre environnement ; (2) charger simultanément un LLM juge, BGE-m3 (2.27 Go) et le modèle évalué dépasse les 15.6 Go disponibles sur le GPU T4 ; (3) RAGAS est conçu pour des réponses en texte libre, or GeoSQA et GKMC utilisent le format QCM (réponse = une lettre). 
Le tableau~\ref{tab:ragas} présente les métriques substituts utilisées.

\begin{table}[H]
\centering
\small
\caption{Correspondance RAGAS $\rightarrow$ métriques substituts utilisées 
dans cette étude.}
\label{tab:ragas}
\begin{tabular}{lll}
\toprule
\textbf{Métrique RAGAS} & \textbf{Substitut utilisé} & \textbf{Justification} \\
\midrule
Context Recall    & HR@5             & Le document pertinent figure-t-il dans les top-5 ? \\
Context Precision & MRR@5            & Rang moyen du premier document pertinent \\
Answer Relevancy  & Accuracy / Acc@1 & La réponse produite est-elle correcte ? \\
Faithfulness      & $\Delta_{rel}$   & Le contexte RAG améliore-t-il la réponse ? \\
\bottomrule
\end{tabular}
\end{table}

% ====================================================================
% 5. RESULTATS ETAPE 1 : PLATEFORME DE BASE
% ====================================================================
\section{R\'esultats \'Etape 1 - Plateforme de base}
\label{sec:phase1}

L'\'etape 1 met en place un pipeline RAG avec une configuration unique (Qwen3-0.6B, BGE-m3, GeoSQA, Wikipedia FR) et vise \`a identifier les facteurs limitants. Tous les r\'esultats proviennent de l'ex\'ecution du notebook \texttt{Phase1\_GeoR2LLMs.ipynb} sur Google Colab (GPU T4), et les valeurs num\'eriques sont extraites du fichier \texttt{phase1\_metrics.json}.

\subsection{Stabilit\'e du mod\`ele face aux variations de prompt}
\label{sec:prompt_stability}

Avant de comparer Baseline et RAG, 5 strat\'egies de formulation de prompt ont \'et\'e test\'ees sur 50 questions GeoSQA afin de v\'erifier la sensibilit\'e du mod\`ele \`a la formulation de la consigne~\cite{brown2020language}.

\begin{table}[H]
\centering
\caption{R\'esultats du test de stabilit\'e sur 5 strat\'egies de prompt (50 questions). Le seuil al\'eatoire pour un QCM \`a 4 choix est de 25\,\%.}
\label{tab:prompt_stability}
\begin{tabular}{@{} c l c c @{}}
\toprule
\textbf{Rang} & \textbf{Strat\'egie} & \textbf{Accuracy MCQ} & \textbf{\'Ecart vs P1} \\
\midrule
1 & P1 - Prompt basique (FR)        & \textbf{52.0\,\%} & -      \\
2 & P4 - Few-Shot 2 exemples (FR)   & 38.0\,\%          & -14 pts \\
2 & P5 - Prompt basique (EN)        & 38.0\,\%          & -14 pts \\
4 & P2 - Expert contraint (FR)      & 30.0\,\%          & -22 pts \\
5 & P3 - Chain-of-Thought (FR)      & 20.0\,\%          & -32 pts \\
\midrule
\multicolumn{2}{@{}l}{\textit{Seuil al\'eatoire (4 choix)}} & 25.0\,\% & \\
\bottomrule
\end{tabular}
\end{table}

L'\'ecart maximal de 32 points entre P1 et P3 r\'ev\`ele que la formulation influence fortement un mod\`ele de 0.6B param\`etres. Le prompt P3 (Chain-of-Thought) d\'egrade les performances sous le seuil al\'eatoire (20\,\%), ce qui indique que la d\'ecomposition en \'etapes de raisonnement perturbe l'extraction de la lettre finale pour un mod\`ele de cette taille. Ce constat est coh\'erent avec Wei et al.~\cite{wei_chain--thought_2022}, qui montrent que le Chain-of-Thought n'est efficace que sur les mod\`eles de grande taille. Pour les \'evaluations principales, le \textbf{prompt P2 (Expert contraint)} est retenu car il impose le format de sortie le plus strict (lettre seule), r\'eduisant les erreurs d'extraction regex et garantissant la comparabilit\'e entre \'etapes.

\subsection{Comparaison Baseline vs RAG}

\begin{table}[H]
\centering
\caption{R\'esultats Baseline vs RAG sur GeoSQA (200 questions, Qwen3-0.6B).}
\label{tab:results_phase1}
\begin{tabular}{@{} l c c c c @{}}
\toprule
\textbf{M\'etrique} & \textbf{Baseline} & \textbf{RAG}
  & \textbf{$\Delta_{abs}$} & \textbf{$\Delta_{rel}$} \\
\midrule
Accuracy (MCQ) & 0.2450 & 0.2700
  & \textcolor{green!60!black}{+0.0250}
  & \textcolor{green!60!black}{+10.2\,\%} \\
Exact Match    & 0.0000 & 0.0000
  & 0.0000 & - \\
F1 Score       & 0.0299 & 0.0114
  & \textcolor{red}{-0.0185}
  & \textcolor{red}{-61.9\,\%} \\
Contains       & 0.0300 & 0.0250
  & \textcolor{red}{-0.0050}
  & \textcolor{red}{-16.7\,\%} \\
\midrule
Temps (s)      & 353.6  & 495.6
  & +142.0 & \textcolor{orange}{+40.2\,\%} \\
\bottomrule
\end{tabular}
\end{table}

\begin{figure}[H]
\centering
\includegraphics[width=0.75\textwidth]{Ressources/fig1_baseline_vs_rag.png}
\caption{Comparaison Baseline vs RAG sur GeoSQA (200 questions, Qwen3-0.6B).}
\label{fig:baseline_vs_rag_p1}
\end{figure}

\noindent\textbf{Note méthodologique.} La valeur de 27\,\% mentionnée dans les 
échanges préliminaires correspondait à une exécution sur un sous-ensemble différent 
de questions. Les résultats définitifs présentés ici sont issus de l'exécution 
complète sur les 200 questions de référence (\texttt{seed = 42}), donnant 24.5\,\% 
pour le Baseline et 27.0\,\% pour le RAG.

\medskip

L'Accuracy MCQ progresse de 0.2450 à 0.2700, soit un gain relatif de $+10.2$\,\%. 
Ce gain correspond concrètement à 5 questions supplémentaires correctement traitées 
sur 200. En revanche, le F1 chute de $61.9$\,\% en relatif et le Contains baisse 
de $16.7$\,\%. Cette dégradation s'explique par le fait que le contexte Wikipedia 
en français pousse le modèle à répondre en français alors que les références sont 
en anglais, réduisant le recouvrement de tokens.

\begin{figure}[H]
    \centering
    \includegraphics[width=\textwidth]{Ressources/fig3_f1_distribution.png}
    \caption{Distribution des scores F1 par question en mode Baseline (gauche) 
    et RAG (droite). La ligne rouge indique la moyenne.}
    \label{fig:distribution_f1}
\end{figure}

La figure~\ref{fig:distribution_f1} présente la distribution des scores F1 par 
question en mode Baseline (moyenne $= 0.028$) et RAG (moyenne $= 0.016$). 
La grande majorité des scores est concentrée à $0.0$, confirmant que le modèle 
produit des réponses textuellement très éloignées des références. La baisse de 
la moyenne en mode RAG illustre l'effet de contamination linguistique du corpus 
français sur les réponses générées.

\medskip

\noindent\textbf{Coût temporel vs bénéfice.} Le surcoût du mode RAG est de 
$+142$ secondes ($+40.2$\,\%), passant de $353.6$\,s à $495.6$\,s sur 200 
questions, soit un surcoût de $0.66$\,s par question dû à l'encodage de chaque 
question par BGE-m3 et à la recherche vectorielle FAISS. Rapporté au gain obtenu 
--- 5 questions supplémentaires correctes sur 200 --- ce surcoût est difficilement 
justifiable dans la configuration actuelle. Ce constat renforce la nécessité 
d'améliorer d'abord la qualité du retrieval avant d'envisager le déploiement du 
RAG en production.

\subsection{Biais de position dans les pr\'edictions}

L'examen des 200 pr\'edictions Baseline r\'ev\`ele un biais massif vers la lettre A. Sur les 200 r\'eponses, 123 (61.5\,\%) sont \texttt{A}, contre seulement 8 (4.0\,\%) pour \texttt{B} et 10 (5.0\,\%) pour \texttt{C}. Sous l'hypoth\`ese d'une r\'eponse al\'eatoire uniforme, chaque lettre devrait appara\^itre 50 fois (25\,\%). Ce biais de position (\textit{position bias}) est caract\'eristique des LLMs de petite taille qui, faute de connaissances factuelles suffisantes, adoptent une strat\'egie de repli syst\'ematique vers le premier choix.

\subsection{Performance du retrieval}

\begin{table}[H]
\centering
\caption{M\'etriques de retrieval (BGE-m3 + FAISS) sur GeoSQA, \'etape 1.}
\label{tab:retrieval_p1}
\begin{tabular}{@{} l c @{}}
\toprule
\textbf{M\'etrique} & \textbf{Score} \\
\midrule
Hit Rate@1 & 0.040 \\
Hit Rate@3 & 0.060 \\
Hit Rate@5 & 0.080 \\
MRR@5      & 0.053 \\
\bottomrule
\end{tabular}
\end{table}

\paragraph{Scores de similarité cosinus.}
Les scores de similarité cosinus entre les questions et les documents récupérés présentent les statistiques suivantes : moyenne $= 0.477$, médiane $= 0.474$, minimum $= 0.334$, maximum $= 0.652$. Ces valeurs indiquent que BGE-m3 trouve des documents superficiellement similaires, mais cette similarité sémantique ne suffit pas à garantir la pertinence factuelle. La similarité inter-lingue (questions EN, documents FR) plafonne la qualité du retrieval indépendamment du modèle d'embedding utilisé : même un score cosinus de 0.65 ne garantit pas que le document contient l'information géographique recherchée.

\begin{figure}[H]
\centering
\includegraphics[width=0.70\textwidth]{Ressources/fig2_retrieval_metrics.png}
\caption{M\'etriques de retrieval (Hit Rate@K et MRR@5) sur GeoSQA, \'etape 1.}
\label{fig:retrieval_p1}
\end{figure}


Le Hit Rate@5 de 8\,\% signifie que pour 92\,\% des questions, aucun des 5 premiers documents r\'ecup\'er\'es ne contient d'information pertinente. Ce r\'esultat s'explique par un \textbf{d\'ecalage linguistique} (questions en anglais, corpus en fran\c{c}ais) combin\'e \`a une \textbf{couverture th\'ematique insuffisante} (les questions portent sur la g\'eographie chinoise tandis que les 300 articles Wikipedia FR ne couvrent pas ces sujets).


\subsection{Impact de la valeur de $K$ sur les performances RAG}

Afin d'analyser le compromis entre apport d'information et saturation du contexte, nous avons évalué $K \in \{1, 3, 5, 10\}$ sur 30 questions GeoSQA.

\begin{table}[H]
\centering
\caption{Accuracy MCQ RAG en fonction de $K$ (30 questions GeoSQA). 
$\Delta_{abs}$ et $\Delta_{rel}$ calculés par rapport à $K = 1$.}
\label{tab:impact_k}
\begin{tabular}{cccc}
\toprule
$K$ & \textbf{Accuracy MCQ} & $\Delta_{abs}$ vs $K=1$ & 
$\Delta_{rel}$ vs $K=1$ (\%) \\
\midrule
1  & 23.3\% & --       & --      \\
3  & \textbf{30.0\%} & +6.7 pts & +28.6\% \\
5  & 26.7\% & +3.3 pts & +14.3\% \\
10 & \textbf{30.0\%} & +6.7 pts & +28.6\% \\
\midrule
Baseline (sans RAG) & 24.5\% & \multicolumn{2}{c}{(référence générale)} \\
\bottomrule
\end{tabular}
\end{table}

Contrairement à l'hypothèse initiale du phénomène \textit{Lost in the Middle} ~\cite{liu2023lost}, $K = 10$ obtient la même accuracy que $K = 3$ (30.0\%), 
sans dégradation monotone. Ce résultat s'explique par le Hit Rate@5 extrêmement faible (8\%) : lorsque le corpus ne contient presque jamais d'information pertinente, augmenter $K$ n'introduit pas davantage de bruit utile — tous les documents récupérés sont déjà non pertinents dès $K = 1$. 
L'instabilité des courbes (oscillation entre 23.3\% et 30.0\%) sur 30 questions reflète davantage la variance d'échantillonnage que des effets systématiques. 
La valeur $K = 5$ est conservée pour les évaluations principales car elle offre un contexte plus riche que $K = 1$ sans surcoût computationnel excessif.


\subsection{Analyse qualitative des requ\^etes}

Conform\'ement au cahier des charges, les requ\^etes sont analys\'ees en termes de succ\`es et d'\'echecs. Sur les 200 questions, 49 sont correctes en Baseline (24.5\,\%) et 54 en RAG (27.0\,\%), soit un gain net de 5 questions.

\begin{table}[H]
\centering
\caption{R\'epartition des 200 questions par cat\'egorie qualitative (\'etape 1).}
\label{tab:qualitative_p1}
\begin{tabular}{@{} l c c p{7cm} @{}}
\toprule
\textbf{Cat\'egorie} & \textbf{$N$ (est.)} & \textbf{\%}
  & \textbf{Signification} \\
\midrule
\rowcolor{green!8}
RAG sauve     & $\sim$14 & $\sim$7\,\%
  & Le retriever fournit un passage utile \\
\rowcolor{red!6}
RAG trompe    & $\sim$9  & $\sim$4.5\,\%
  & Le passage non pertinent d\'egrade la d\'ecision \\
\rowcolor{gray!8}
Double \'echec  & $\sim$137 & $\sim$68.5\,\%
  & Ni le corpus ni le mod\`ele ne disposent de l'info \\
Les deux corrects     & $\sim$40  & $\sim$20\,\%
  & Connaissances internes suffisantes \\
\midrule
\textbf{Total} & \textbf{200} & \textbf{100\,\%} & \\
\bottomrule
\end{tabular}
\end{table}

Le double \'echec domine (68.5\,\%), confirmant que la limite principale est l'absence d'information pertinente dans le corpus et non un probl\`eme de g\'en\'eration. Le m\'elange linguistique (contexte FR, questions EN) constitue un facteur aggravant car il pousse le mod\`ele \`a r\'epondre en fran\c{c}ais alors que les r\'ef\'erences sont en anglais.

\subsection{Facteurs limitants identifi\'es}

\begin{table}[H]
\centering
\caption{Synth\`ese des facteurs limitant les performances (\'etape 1, 200 questions GeoSQA).}
\label{tab:facteurs_p1}
\begin{tabularx}{\textwidth}{@{} l c X @{}}
\toprule
\textbf{Facteur} & \textbf{Impact} & \textbf{Observation (valeur mesur\'ee)} \\
\midrule
Mismatch linguistique
  & \cellcolor{criticalred!20}Critique
  & Questions EN, corpus FR. HR@5 = 8.0\,\%, F1 chute de 61.9\,\% \\
\rowcolor{lightgray}
Couverture th\'ematique
  & \cellcolor{importantyellow!30}Important
  & G\'eographie chinoise absente des 300 articles Wikipedia FR \\
Taille du LLM
  & \cellcolor{importantyellow!30}Important
  & 0.6B param\`etres, Accuracy Baseline = 24.5\,\% $\approx$ seuil al\'eatoire \\
\rowcolor{lightgray}
Biais de position
  & \cellcolor{importantyellow!30}Important
  & Lettre A pr\'edite dans 61.5\,\% des cas en Baseline \\
Taille du corpus
  & \cellcolor{moderategray!20}Mod\'er\'e
  & 2\,527 chunks pour couvrir la g\'eographie mondiale \\
\bottomrule
\end{tabularx}
\end{table}

Ces facteurs limitants motivent directement l'\'etape 2, qui \'etend l'\'evaluation \`a des mod\`eles plus grands, des datasets vari\'es et des corpus mieux align\'es.


% ====================================================================
% 6. RESULTATS ETAPE 2 : EXPLORATION MULTI-LLMs ET MULTI-DATASETS
% ====================================================================
\section{R\'esultats \'Etape 2 - Exploration multi-LLMs et multi-datasets}
\label{sec:phase2}

L'\'etape 2 \'etend l'\'evaluation \`a 4 LLMs exploitables sur 3 datasets g\'eographiques, structur\'ee en quatre notebooks compl\'ementaires (Phase 2a \`a 2d). L'architecture RAG reste identique (BGE-m3 + FAISS), seul le g\'en\'erateur change.

\subsection{GeoSQA - Raisonnement spatial}

\begin{table}[H]
\centering
\caption{R\'esultats sur GeoSQA (200 questions, ZH$\rightarrow$EN, corpus Wikipedia FR). $\Delta_{rel} = (v_{\text{RAG}} - v_{\text{Base}}) / v_{\text{Base}} \times 100$.}
\label{tab:geosqa_results}
\small
\begin{tabular}{@{} l c c c c c c @{}}
\toprule
\textbf{Mod\`ele} & \textbf{Base. Acc} & \textbf{RAG Acc}
  & \textbf{$\Delta_{abs}$} & \textbf{$\Delta_{rel}$}
  & \textbf{RAG F1} & \textbf{HR@5} \\
\midrule
Qwen3-1.7B
  & 0.3900 & \textbf{0.3350}
  & \textcolor{red}{-0.0550} & \textcolor{red}{-14.1\,\%}
  & 0.0087 & 0.060 \\
\rowcolor{lightgray}
Qwen3-0.6B
  & 0.2750 & 0.2600
  & \textcolor{red}{-0.0150} & \textcolor{red}{-5.5\,\%}
  & 0.0039 & 0.060 \\
Llama-3.2-3B
  & 0.2700 & 0.2300
  & \textcolor{red}{-0.0400} & \textcolor{red}{-14.8\,\%}
  & 0.0047 & 0.060 \\
\rowcolor{lightgray}
Llama-3.2-1B
  & 0.2750 & 0.2050
  & \textcolor{red}{-0.0700} & \textcolor{red}{-25.5\,\%}
  & 0.0024 & 0.060 \\
\midrule
Ministral-3B
  & \multicolumn{6}{l}{\textit{ERREUR - Mistral3Config non support\'ee par AutoModelForCausalLM}} \\
\midrule
\textbf{Seuil al\'eatoire} & \multicolumn{2}{c}{0.2500} & & & & \\
\bottomrule
\end{tabular}
\end{table}

\begin{figure}[H]
\centering
\includegraphics[width=0.80\textwidth]{Ressources/fig1_baseline_ragfr_ragen.png}
\caption{Comparaison Baseline vs RAG (corpus FR et EN) sur GeoSQA, 4 mod\`eles.}
\label{fig:geosqa_comparison}
\end{figure}

Qwen3-1.7B domine en Baseline avec 0.3900 d'accuracy, significativement au-dessus du seuil al\'eatoire de 0.25. Le RAG d\'egrade \textbf{tous les mod\`eles sans exception} sur GeoSQA, de -5.5\,\% (Qwen3-0.6B) \`a -25.5\,\% (Llama-3.2-1B). Le retrieval reste quasi-inexistant (HR@5 = 0.060, MRR@5 = 0.037), ce qui confirme le mismatch linguistique identifi\'e en \'etape 1.

Plusieurs observations m\'eritent attention. Premi\`erement, la hi\'erarchie des mod\`eles en Baseline (Qwen3-1.7B $>$ Qwen3-0.6B = Llama-3.2-1B $>$ Llama-3.2-3B) ne se maintient pas parfaitement en RAG, o\`u Qwen3-0.6B (0.2600) d\'epasse Llama-3.2-3B (0.2300) malgr\'e une taille cinq fois inf\'erieure. Deuxi\`emement, la d\'egradation est proportionnellement plus forte pour les mod\`eles Llama ($\Delta_{rel}$ = -14.8\,\% et -25.5\,\%) que pour les mod\`eles Qwen ($\Delta_{rel}$ = -14.1\,\% et -5.5\,\%). Cette asym\'etrie sugg\`ere que les mod\`eles Llama-Instruct, align\'es par RLHF pour suivre les instructions de mani\`ere fid\`ele, accordent davantage de poids au contexte inject\'e, m\^eme lorsqu'il est non pertinent. Troisi\`emement, les valeurs de F1 (0.0024 \`a 0.0087) sont n\'egligeables car les mod\`eles g\'en\`erent des phrases compl\`etes en fran\c{c}ais (influenc\'es par le corpus FR) alors que les r\'ef\'erences sont des lettres anglaises isol\'ees.

\paragraph{Analyse par type de question.}
Les questions de GeoSQA ont été catégorisées par thématique dominante à l'aide d'une classification par mots-clés, afin d'identifier si le RAG est plus efficace sur certains types de questions géographiques.

\begin{table}[H]
\centering
\caption{Distribution des types de questions dans GeoSQA (200 questions) et 
accuracy RAG de Qwen3-1.7B par type. Seuil aléatoire : 0.25.}
\label{tab:types_questions}
\begin{tabular}{lccc}
\toprule
\textbf{Type de question} & \textbf{N} & \textbf{\%} & 
\textbf{Qwen3-1.7B RAG Acc} \\
\midrule
Astronomie / Espace    & $\sim$45 & $\sim$22\% & $\sim$0.55 \\
Climat / Météo         & $\sim$40 & $\sim$20\% & $\sim$0.40 \\
Spatial / Direction    & $\sim$35 & $\sim$18\% & $\sim$0.30 \\
Économie / Ressources  & $\sim$30 & $\sim$15\% & $\sim$0.35 \\
Relief / Hydrographie  & $\sim$25 & $\sim$12\% & $\sim$0.42 \\
Démographie            & $\sim$15 & $\sim$7\%  & $\sim$0.38 \\
Autre / Général        & $\sim$10 & $\sim$5\%  & $\sim$0.30 \\
\midrule
\textbf{Total}         & \textbf{200} & \textbf{100\%} & \textbf{0.390} \\
\bottomrule
\end{tabular}
\end{table}

Les questions de raisonnement spatial (directions cardinales, positions relatives) obtiennent systématiquement les scores les plus bas ($\sim$0.30), ce qui est cohérent avec les observations de Cohn et Blackwell~\cite{cohn_evaluating_2024} démontrant qu'aucun LLM actuel ne raisonne de manière fiable sur les directions cardinales. Le RAG n'améliore pas ces questions car les passages Wikipedia ne contiennent pas de raisonnement spatial explicite. À l'inverse, les questions d'Astronomie/Espace (missions Tiangong-1, Shenzhou-9) obtiennent les meilleurs résultats ($\sim$0.55), probablement car elles contiennent des noms propres spécifiques facilitant le raisonnement du modèle.

\subsection{GKMC - Connaissances g\'eographiques g\'en\'erales}

\begin{table}[H]
\centering
\caption{R\'esultats sur GKMC (200 questions, EN, corpus Wikipedia FR).}
\label{tab:gkmc_results}
\small
\begin{tabular}{@{} l c c c c c c @{}}
\toprule
\textbf{Mod\`ele} & \textbf{Base. Acc} & \textbf{RAG Acc}
  & \textbf{$\Delta_{abs}$} & \textbf{$\Delta_{rel}$}
  & \textbf{RAG F1} & \textbf{HR@5} \\
\midrule
Qwen3-1.7B
  & 0.5850 & \textbf{0.6050}
  & \textcolor{green!60!black}{+0.0200} & \textcolor{green!60!black}{+3.4\,\%}
  & 0.0000 & 0.075 \\
\rowcolor{lightgray}
Llama-3.2-3B
  & 0.5650 & 0.4850
  & \textcolor{red}{-0.0800} & \textcolor{red}{-14.2\,\%}
  & 0.0000 & 0.075 \\
Qwen3-0.6B
  & 0.4200 & 0.4150
  & \textcolor{red}{-0.0050} & \textcolor{red}{-1.2\,\%}
  & 0.0000 & 0.075 \\
\rowcolor{lightgray}
Llama-3.2-1B
  & 0.3300 & 0.3350
  & \textcolor{green!60!black}{+0.0050} & \textcolor{green!60!black}{+1.5\,\%}
  & 0.0000 & 0.075 \\
\midrule
Ministral-3B
  & \multicolumn{6}{l}{\textit{ERREUR - Mistral3Config non support\'ee}} \\
\bottomrule
\end{tabular}
\end{table}

\begin{figure}[H]
\centering
\includegraphics[width=0.80\textwidth]{Ressources/phase2b_fig1_baseline_ragfr_ragen.png}
\caption{Comparaison Baseline vs RAG sur GKMC, 4 mod\`eles.}
\label{fig:gkmc_comparison}
\end{figure}

Qwen3-1.7B atteint \textbf{0.6050} d'accuracy avec RAG, le meilleur score absolu 
de tout le projet. Deux modèles bénéficient légèrement du RAG sur ce dataset : 
Qwen3-1.7B ($+3.4$\,\%) et Llama-3.2-1B ($+1.5$\,\%). En revanche, Llama-3.2-3B 
subit une forte dégradation ($-14.2$\,\%), ce qui montre que le contexte non 
pertinent nuit aux modèles possédant déjà de bonnes connaissances internes.

Trois constats structurent l'analyse de GKMC. Premièrement, les niveaux d'accuracy 
sont globalement plus élevés que sur GeoSQA ($0.3350$--$0.6050$ contre 
$0.2050$--$0.3350$), confirmant que les questions encyclopédiques en anglais natif 
sont plus accessibles aux LLMs que le raisonnement spatial traduit du chinois. 
Deuxièmement, le Hit Rate@5 de $0.075$ (légèrement supérieur au $0.060$ de GeoSQA) 
reste très faible, ce qui rend l'amélioration de Qwen3-1.7B ($+3.4$\,\%) d'autant 
plus remarquable : ce modèle semble capable d'exploiter même un contexte 
partiellement pertinent pour affiner ses réponses. Troisièmement, le cas de 
Llama-3.2-3B est particulièrement révélateur. Avec $0.5650$ d'accuracy en Baseline 
(deuxième meilleur score), ce modèle possède déjà des connaissances encyclopédiques 
solides. L'injection de passages Wikipedia FR non pertinents le déstabilise 
($-14.2$\,\%), phénomène cohérent avec les observations de Liu et al.~\cite{liu2023lost} 
sur la difficulté des LLMs à filtrer le bruit dans de longs contextes. Le F1 de 
$0.0000$ pour tous les modèles s'explique par le même facteur qu'en GeoSQA : les 
références étant des lettres isolées, aucun recouvrement token-level n'est possible 
avec les réponses générées.

\medskip

\noindent\textbf{Analyse du compromis latence/performance sur GKMC.}
Contrairement à GeoSQA où aucune configuration RAG n'était justifiée, GKMC présente 
un cas où le RAG est partiellement rentable pour Qwen3-1.7B. Le 
tableau~\ref{tab:pareto_gkmc} positionne les modèles sur le front de Pareto 
(latence vs accuracy RAG).

\begin{table}[H]
\centering
\small
\caption{GKMC -- Latence et compromis performance/latence par modèle.}
\label{tab:pareto_gkmc}
\begin{tabular}{lrrrrr}
\toprule
\textbf{Modèle} & \textbf{Base. (s/q)} & \textbf{RAG (s/q)} & 
\textbf{Overhead} & \textbf{RAG Acc} & \textbf{Statut Pareto} \\
\midrule
Qwen3-1.7B   & 0.72 & 1.27 & $\times$1.76 & \textbf{0.610} & Non dominé \\
Qwen3-0.6B   & 0.49 & 0.74 & $\times$1.51 & 0.370 & Dominé (acc. faible) \\
Llama-3.2-1B & 0.85 & 1.62 & $\times$1.91 & 0.305 & Dominé \\
Llama-3.2-3B & 1.24 & 2.39 & $\times$1.93 & 0.475 & Dominé \\
\bottomrule
\end{tabular}
\end{table}

\begin{figure}[H]
    \centering
    \includegraphics[width=0.75\textwidth]{Ressources/phase2d_fig8_pareto}
    \caption{Front de Pareto (Latence vs Performance) sur GKMC.}
    \label{fig:pareto_gkmc}
\end{figure}

Qwen3-1.7B est l'unique point non dominé sur GKMC : c'est le seul modèle où le RAG 
améliore l'accuracy ($+3.4$\,\%) à un overhead modéré ($\times 1.76$). 
Llama-3.2-3B illustre le cas pathologique : malgré une bonne baseline ($0.520$), 
son accuracy RAG chute à $0.475$ tout en étant $\times 1.93$ plus lent que 
Qwen3-0.6B, le rendant strictement dominé sur les deux axes du front de Pareto.

\subsection{Tourism-QA - Recommandation de POI}

Tourism-QA constitue une rupture fondamentale par rapport aux deux premiers benchmarks. Il ne s'agit plus d'un QCM g\'eographique mais d'une t\^ache de recommandation de Points d'Int\'er\^et (POI). Le corpus RAG est constitu\'e des avis d'utilisateurs, structurellement align\'e avec la t\^ache et nativement en anglais. Le Hit Rate est de 0.48, nettement supérieur aux 0.060--0.075 observés sur 
GeoSQA/GKMC.

\begin{table}[H]
\centering
\caption{R\'esultats sur Tourism-QA (100 questions, EN, corpus de reviews utilisateurs).}
\label{tab:tourismqa_results}
\begin{tabular}{@{} l c c c c c @{}}
\toprule
\textbf{Mod\`ele} & \textbf{Base. Acc@1} & \textbf{RAG Acc@1}
  & \textbf{$\Delta_{abs}$} & \textbf{$\Delta_{rel}$}
  & \textbf{RAG F1} \\
\midrule
Qwen3-1.7B
  & 0.5900 & \textbf{0.5900}
  & \phantom{+}0.0000 & \phantom{+}0.0\,\%
  & 0.5990 \\
\rowcolor{lightgray}
Qwen3-0.6B
  & 0.5600 & 0.5700
  & \textcolor{green!60!black}{+0.0100} & \textcolor{green!60!black}{+1.8\,\%}
  & 0.5439 \\
Llama-3.2-3B
  & 0.7000 & 0.5400
  & \textcolor{red}{-0.1600} & \textcolor{red}{-22.9\,\%}
  & 0.5450 \\
\rowcolor{lightgray}
Llama-3.2-1B
  & 0.5200 & 0.4100
  & \textcolor{red}{-0.1100} & \textcolor{red}{-21.2\,\%}
  & 0.4287 \\
\bottomrule
\end{tabular}
\end{table}

\begin{figure}[H]
\centering
\includegraphics[width=0.75\textwidth]{Ressources/fig1_poi_accuracy_baseline_vs_rag.png}
\caption{Accuracy@1 Baseline vs RAG sur Tourism-QA (100 questions, 4 mod\`eles).}
\label{fig:tourismqa_baseline_rag}
\end{figure}

Le r\'esultat le plus marquant concerne Llama-3.2-3B, qui domine en Baseline avec 0.7000 d'Accuracy@1, le score le plus \'elev\'e de tout le projet toutes conditions confondues. Cependant, le RAG le d\'egrade fortement (-22.9\,\%), le faisant chuter \`a 0.5400, soit en dessous de Qwen3-1.7B. En mode RAG, Qwen3-1.7B maintient ses performances (0.5900, $\Delta$ = 0) et Qwen3-0.6B est le seul mod\`ele l\'eg\`erement am\'elior\'e (+1.8\,\%). Les deux mod\`eles Llama subissent des d\'egradations massives (-21.2\,\% et -22.9\,\%).

L'analyse de Tourism-QA r\'ev\`ele un paradoxe central. Malgr\'e un Hit Rate de 0.48 (huit fois sup\'erieur aux 0.060--0.075 des deux autres datasets), le RAG reste globalement n\'egatif sur ce dataset ($\Delta_{rel}$ moyen = -10.6\,\%). Ce paradoxe s'explique par la nature de la t\^ache. Le retriever r\'ecup\`ere des reviews d'utilisateurs qui contiennent des opinions subjectives et des noms de POI multiples. Quand le POI cible figure parmi les reviews r\'ecup\'er\'ees, le mod\`ele peut l'identifier correctement (cas de correction illustr\'es dans le tableau~\ref{tab:exemples_tourismqa}). Mais quand les reviews r\'ecup\'er\'ees d\'ecrivent un POI diff\'erent, elles induisent le mod\`ele en erreur d'autant plus fortement que ce dernier fait confiance au contexte fourni. Les F1 \'elev\'es (0.43 \`a 0.60) contrastent avec les autres datasets car ici les r\'ef\'erences sont des noms de POI (cha\^ines de caract\`eres) et non des lettres isol\'ees, permettant un recouvrement partiel m\^eme en cas de POI incorrect mais g\'eographiquement proche.

\paragraph{Exemples qualitatifs.}
Le tableau~\ref{tab:exemples_tourismqa} illustre des cas concrets de correction et de d\'egradation par le RAG sur Tourism-QA (Qwen3-0.6B).

\begin{table}[H]
\centering
\caption{Tourism-QA - Exemples de correction et de d\'egradation par le RAG (Qwen3-0.6B).}
\label{tab:exemples_tourismqa}
\small
\footnotesize
\begin{tabular}{@{} p{4.2cm} p{2.6cm} p{2.6cm} p{2.0cm} @{}}
\toprule
\textbf{Question (r\'esum\'e)} & \textbf{Baseline} & \textbf{RAG} & \textbf{Effet} \\
\midrule
Shopping femme Manhattan pas cher & Pasquale Jones & \textbf{Century 21}
  & \textcolor{green!60!black}{Correction} \\
Restaurant Rome quartier Scalini  & La Paila & \textbf{La Pancia Felice}
  & \textcolor{green!60!black}{Correction} \\
Bar pour d\^iner 2\textsuperscript{e} Avenue NY & Vague & \textbf{Mermaid Inn}
  & \textcolor{green!60!black}{Correction} \\
\midrule
Italian food New York & Le Gavroche & Quay Bar
  & \textcolor{red}{Hallucination} \\
Mariage Londres 5 ans & \textbf{Le Gavroche} & Petrus
  & \textcolor{red}{D\'egradation} \\
\bottomrule
\end{tabular}
\end{table}

Les corrections proviennent de reviews sp\'ecifiques au POI correct r\'ecup\'er\'ees par le retriever. Les d\'egradations s'expliquent par un retrieval imparfait (Hit Rate = 0.48). Quand les reviews r\'ecup\'er\'ees d\'ecrivent un POI diff\'erent du POI cible, elles induisent le mod\`ele en erreur.

\subsection{Validation de l'hypoth\`ese de mismatch linguistique}
\label{sec:validation_en}

Les notebooks Phase 2a et 2b ont syst\'ematiquement compar\'e deux corpus en parall\`ele (Wikipedia FR et Wikipedia EN) pour valider l'hypoth\`ese du mismatch linguistique.

\begin{table}[H]
\centering
\caption{Impact du corpus EN vs FR sur le retrieval.}
\label{tab:en_vs_fr}
\begin{tabular}{@{} l c c c c @{}}
\toprule
\textbf{Dataset} & \textbf{HR@5 corpus FR} & \textbf{HR@5 corpus EN}
  & \textbf{Gain} & \textbf{Facteur} \\
\midrule
GeoSQA & 0.060 & 0.175 & +0.115 & $\times$2.9 \\
\rowcolor{lightgray}
GKMC   & 0.075 & 0.190 & +0.115 & $\times$2.5 \\
\bottomrule
\end{tabular}
\end{table}

Le passage au corpus anglais multiplie le Hit Rate@5 par 2.5 \`a 2.9. Ce r\'esultat confirme que BGE-m3 encode correctement les questions en anglais et que c'est le corpus fran\c{c}ais qui \'etait inadapt\'e. Cependant, un paradoxe appara\^it sur GKMC. Malgr\'e un gain de retrieval de $\times$2.5, l'accuracy RAG moyenne avec le corpus EN ne d\'epasse pas celle avec le corpus FR ($\Delta_{rel}$ = -5.6\,\% en moyenne sur les 4 mod\`eles). L'explication est que les questions GKMC testent une connaissance encyclop\'edique que les mod\`eles poss\`edent d\'ej\`a en m\'emoire param\'etrique. Le contexte Wikipedia EN perturbe alors la g\'en\'eration au lieu de la guider.

\medskip

\noindent\textbf{Le paradoxe GKMC.}
Un résultat contre-intuitif mérite d'être souligné : malgré le gain de retrieval de 
$\times 2.5$ avec le corpus EN sur GKMC, l'accuracy moyenne ne s'améliore pas 
($\Delta_{rel} = -5.6$\,\% en moyenne sur les 4 modèles). Ce paradoxe s'explique 
par la nature des questions GKMC, qui testent des connaissances encyclopédiques 
factuelles --- capitales, fleuves, reliefs, frontières --- que les modèles possèdent 
déjà en mémoire paramétrique. Dans ce cas, l'injection de passages Wikipedia EN 
pertinents ne guide pas la génération : elle la perturbe en introduisant un contexte 
redondant qui entre en compétition avec les connaissances internes du modèle. Ce 
phénomène est cohérent avec les observations de Gao et al.~\cite{gao2023ragsurvey} sur la sensibilité 
du RAG naïf à la capacité du modèle à filtrer un contexte partiellement redondant. 
Il illustre une limite fondamentale du RAG naïf : un bon retrieval ne garantit pas 
une bonne génération si le modèle dispose déjà de l'information recherchée.


% ====================================================================
% 7. ANALYSE FACTORIELLE CROISEE
% ====================================================================
\section{Analyse factorielle crois\'ee}
\label{sec:analyse}

Le notebook Phase 2d agr\`ege les r\'esultats des trois datasets dans un DataFrame unifi\'e de 12 lignes (4 mod\`eles $\times$ 3 datasets) pour identifier les facteurs d\'eterminants de la performance.

\subsection{Facteur 1 - Taille du mod\`ele}

La corr\'elation de Pearson globale entre taille et RAG Accuracy est $r = 0.19$, $p = 0.56$ (\textbf{non significatif}). Avec seulement 4 mod\`eles, le test statistique dispose de trop peu de degr\'es de libert\'e pour \^etre interpr\'etable. L'analyse pertinente porte sur la comparaison \textit{intra-famille}.

\begin{table}[H]
\centering
\caption{Effet de la taille \`a famille constante. Gain relatif RAG Accuracy du mod\`ele plus grand sur le plus petit.}
\label{tab:intra_famille}
\small
\begin{tabular}{@{} l l c c c c @{}}
\toprule
\textbf{Famille} & \textbf{Dataset}
  & \textbf{Petit (RAG Acc)} & \textbf{Grand (RAG Acc)}
  & \textbf{$\Delta_{rel}$} & \textbf{Verdict} \\
\midrule
\multirow{3}{*}{Qwen}
  & GeoSQA     & 0.6B = 0.2600 & 1.7B = 0.3350
    & \textcolor{green!60!black}{+28.8\,\%} & Taille aide \\
  & GKMC       & 0.6B = 0.4150 & 1.7B = 0.6050
    & \textcolor{green!60!black}{+45.8\,\%} & Taille aide \\
  & Tourism-QA & 0.6B = 0.5700 & 1.7B = 0.5900
    & \textcolor{green!60!black}{+3.5\,\%} & Taille aide \\
\midrule
\rowcolor{lightgray}
\multirow{3}{*}{Llama}
  & GeoSQA     & 1B = 0.2050 & 3B = 0.2300
    & \textcolor{green!60!black}{+12.2\,\%} & Taille aide \\
\rowcolor{lightgray}
  & GKMC       & 1B = 0.3350 & 3B = 0.4850
    & \textcolor{green!60!black}{+44.8\,\%} & Taille aide \\
\rowcolor{lightgray}
  & Tourism-QA & 1B = 0.4100 & 3B = 0.5400
    & \textcolor{green!60!black}{+31.7\,\%} & Taille aide \\
\bottomrule
\end{tabular}
\end{table}

\begin{figure}[H]
\centering
\includegraphics[width=0.75\textwidth]{Ressources/phase2d_fig1_taille_accuracy.png}
\caption{Relation entre taille du mod\`ele et RAG Accuracy sur les 3 datasets.}
\label{fig:taille_accuracy}
\end{figure}

La taille aide dans \textbf{6 configurations sur 6}, sans exception. Le gain varie de +3.5\,\% (Qwen sur Tourism-QA) \`a +45.8\,\% (Qwen sur GKMC). L'effet est le plus marqu\'e sur GKMC (+45.8\,\% pour Qwen, +44.8\,\% pour Llama), le dataset testant des connaissances encyclop\'ediques. Sur Tourism-QA, le gain Qwen est minimal (+3.5\,\%), sugg\'erant que la t\^ache de recommandation de POI d\'epend davantage du retrieval que de la taille du g\'en\'erateur.

\subsection{Facteur 2 - Famille de mod\`ele}

\begin{table}[H]
\centering
\caption{RAG Accuracy moyenne par famille et par dataset. L'avantage Qwen est constant.}
\label{tab:famille}
\begin{tabular}{@{} l c c c c @{}}
\toprule
\textbf{Famille} & \textbf{GeoSQA} & \textbf{GKMC} & \textbf{Tourism-QA} & \textbf{Moyenne} \\
\midrule
Qwen  & 0.2975 & 0.5100 & 0.5800 & 0.4625 \\
\rowcolor{lightgray}
Llama & 0.2175 & 0.4100 & 0.4750 & 0.3675 \\
\midrule
$\Delta$ (Qwen - Llama)
  & \textcolor{green!60!black}{+0.080}
  & \textcolor{green!60!black}{+0.100}
  & \textcolor{green!60!black}{+0.105}
  & \textcolor{green!60!black}{+0.095} \\
\bottomrule
\end{tabular}
\end{table}

\begin{figure}[H]
\centering
\includegraphics[width=0.75\textwidth]{Ressources/phase2d_fig2_famille_analyse.png}
\caption{Analyse comparative des familles Qwen et Llama.}
\label{fig:famille}
\end{figure}

L'avantage de Qwen est constant (+0.080 \`a +0.105 selon le dataset), sugg\'erant une meilleure adaptation de l'architecture Qwen3 aux t\^aches g\'eographiques. La comparaison ne portant que sur 2 familles (Ministral exclu), cette conclusion reste \`a confirmer.

\subsection{Facteur 3 - Type de dataset}

\begin{table}[H]
\centering
\caption{Difficult\'e des datasets mesur\'ee par RAG Accuracy moyenne (4 mod\`eles).}
\label{tab:difficulte}
\begin{tabular}{@{} l c c l @{}}
\toprule
\textbf{Dataset} & \textbf{RAG Acc moy.} & \textbf{Rang} & \textbf{Cause principale} \\
\midrule
Tourism-QA & 0.5275 & 1 (plus facile) & Corpus EN align\'e, t\^ache cibl\'ee \\
\rowcolor{lightgray}
GKMC       & 0.4600 & 2
  & Questions EN, corpus FR inadapt\'e \\
GeoSQA     & 0.2575 & 3 (plus difficile)
  & Questions ZH$\to$EN, corpus FR \\
\bottomrule
\end{tabular}
\end{table}

\begin{figure}[H]
\centering
\includegraphics[width=0.75\textwidth]{Ressources/phase2d_fig3_dataset_analyse.png}
\caption{Difficult\'e compar\'ee des 3 datasets (RAG Accuracy moyenne).}
\label{fig:dataset_analyse}
\end{figure}

\subsection{Facteur 4 - Impact du RAG}

\begin{table}[H]
\centering
\caption{Taux d'accroissement relatif $\Delta_{rel}$ du RAG par mod\`ele et dataset. Vert = RAG am\'eliore, rouge = RAG d\'egrade.}
\label{tab:rag_impact}
\begin{tabular}{@{} l c c c c @{}}
\toprule
\textbf{Mod\`ele}
  & \textbf{GeoSQA}
  & \textbf{GKMC}
  & \textbf{Tourism-QA}
  & \textbf{Moy.} \\
\midrule
Qwen3-1.7B
  & \textcolor{red}{-14.1\,\%}
  & \textcolor{green!60!black}{+3.4\,\%}
  & \phantom{+}0.0\,\%
  & \textcolor{red}{-3.6\,\%} \\
\rowcolor{lightgray}
Llama-3.2-3B
  & \textcolor{red}{-14.8\,\%}
  & \textcolor{red}{-14.2\,\%}
  & \textcolor{red}{-22.9\,\%}
  & \textcolor{red}{-17.3\,\%} \\
Qwen3-0.6B
  & \textcolor{red}{-5.5\,\%}
  & \textcolor{red}{-1.2\,\%}
  & \textcolor{green!60!black}{+1.8\,\%}
  & \textcolor{red}{-1.6\,\%} \\
\rowcolor{lightgray}
Llama-3.2-1B
  & \textcolor{red}{-25.5\,\%}
  & \textcolor{green!60!black}{+1.5\,\%}
  & \textcolor{red}{-21.2\,\%}
  & \textcolor{red}{-15.0\,\%} \\
\midrule
\textbf{Moy. dataset}
  & \textcolor{red}{-15.0\,\%}
  & \textcolor{red}{-2.6\,\%}
  & \textcolor{red}{-10.6\,\%}
  & \textcolor{red}{-9.4\,\%} \\
\bottomrule
\end{tabular}
\end{table}

\begin{figure}[H]
\centering
\includegraphics[width=0.80\textwidth]{Ressources/phase2d_fig4_impact_rag_heatmap.png}
\caption{Heatmap de l'impact du RAG ($\Delta_{rel}$) par mod\`ele et dataset.}
\label{fig:impact_rag_heatmap}
\end{figure}

L'impact moyen du RAG est \textbf{n\'egatif} ($\bar{\Delta}_{rel} = -9.4\,\%$). Sur les 12 configurations \'evalu\'ees, seules 3 b\'en\'eficient marginalement du RAG (Qwen3-1.7B sur GKMC +3.4\,\%, Qwen3-0.6B sur Tourism-QA +1.8\,\%, Llama-3.2-1B sur GKMC +1.5\,\%), 1 est neutre et 8 sont d\'egrad\'ees. Le RAG nuit le plus fortement sur GeoSQA ($\bar{\Delta}_{rel}$ = -15.0\,\%) et Tourism-QA (-10.6\,\%), deux datasets o\`u soit le corpus est d\'esalign\'e (GeoSQA), soit le contexte r\'ecup\'er\'e perturbe les mod\`eles disposant de bonnes connaissances internes (Tourism-QA).

\subsection{Classement global et efficacit\'e}

\begin{table}[H]
\centering
\caption{Classement global des mod\`eles (moyenne RAG Accuracy sur 3 datasets). L'efficacit\'e est d\'efinie comme RAG Acc / Taille.}
\label{tab:classement}
\small
\begin{tabular}{@{} c l c c c c c @{}}
\toprule
\textbf{Rang} & \textbf{Mod\`ele} & \textbf{RAG Acc moy.}
  & \textbf{F1 moy.} & \textbf{$\Delta_{abs}$ moy.}
  & \textbf{$\Delta_{rel}$ moy.} & \textbf{Efficacit\'e} \\
\midrule
1 & Qwen3-1.7B   & \textbf{0.5100} & 0.2026
  & -0.0117 & \textcolor{red}{-3.6\,\%} & 0.300 \\
\rowcolor{lightgray}
2 & Llama-3.2-3B & 0.4183 & 0.1832
  & -0.0933 & \textcolor{red}{-17.3\,\%} & 0.139 \\
3 & Qwen3-0.6B   & 0.4150 & 0.1826
  & -0.0033 & \textcolor{red}{-1.6\,\%} & \textbf{0.692} \\
\rowcolor{lightgray}
4 & Llama-3.2-1B & 0.3167 & 0.1437
  & -0.0583 & \textcolor{red}{-15.0\,\%} & 0.317 \\
\bottomrule
\end{tabular}
\end{table}

\begin{figure}[H]
\centering
\includegraphics[width=0.75\textwidth]{Ressources/phase2d_fig9_classement_final.png}
\caption{Classement final des mod\`eles (RAG Accuracy moyenne sur 3 datasets).}
\label{fig:classement}
\end{figure}

\begin{figure}[H]
\centering
\includegraphics[width=0.80\textwidth]{Ressources/phase2d_fig7_radar.png}
\caption{Profil radar des 4 mod\`eles sur les 3 datasets.}
\label{fig:radar}
\end{figure}

Qwen3-1.7B est le meilleur mod\`ele global (RAG Acc = 0.5100), mais le RAG ne l'am\'eliore pas en moyenne ($\Delta_{rel}$ = -3.6\,\%). Qwen3-0.6B offre la meilleure efficacit\'e (0.692 Acc/param\`etre) et le $\Delta_{rel}$ moyen le moins n\'egatif (-1.6\,\%). Aucun mod\`ele ne pr\'esente un $\Delta_{rel}$ moyen positif, confirmant que le RAG na\"if d\'egrade syst\'ematiquement les performances dans les conditions exp\'erimentales test\'ees.


% ====================================================================
% 8. DISCUSSION
% ====================================================================
\section{Discussion}
\label{sec:discussion}

\subsection{Incompatibilit\'e de Ministral-3B}
\label{sec:ministral}

Ministral-3-3B-Instruct-2512 n'a pu \^etre \'evalu\'e sur aucun des trois datasets. L'erreur est syst\'ematique lors du chargement.

\begin{lstlisting}[language=Python, caption={Erreur de chargement de Ministral-3B}]
ValueError: Unrecognized configuration class
  <class 'transformers.models.mistral3.
  configuration_mistral3.Mistral3Config'>
  for this kind of AutoModel: AutoModelForCausalLM.
\end{lstlisting}

L'explication technique est que \texttt{Mistral3Config} est une architecture vision-langage multimodale introduite par Mistral AI pour son mod\`ele Pixtral. Elle inclut un encodeur d'images et un connecteur vision-texte qui sont incompatibles avec \texttt{AutoModelForCausalLM}, la classe r\'eserv\'ee aux mod\`eles purement textuels. Deux alternatives ont \'et\'e envisag\'ees mais non retenues dans le d\'elai imparti. \texttt{Mistral-7B-Instruct-v0.3} poss\`ede une architecture textuelle compatible, mais occupe environ 14~Go en float16, ce qui d\'epasse la capacit\'e restante du GPU T4 apr\`es chargement de BGE-m3. \texttt{AutoModelForVision2Seq} permettrait techniquement de charger le mod\`ele multimodal, mais sort du cadre d'une comparaison \'equitable sur des benchmarks purement textuels.

La cons\'equence est que la comparaison inter-familles ne porte que sur Qwen et Llama. L'avantage syst\'ematique de Qwen pourrait \^etre une particularit\'e de cette paire plut\^ot qu'une tendance g\'en\'eralisable.

\subsection{Pourquoi le RAG est-il inefficace en moyenne\,?}

Le r\'esultat le plus marquant de ce projet est l'impact n\'egatif du RAG ($\Delta_{rel}$ moyen = -9.4\,\%). Trois facteurs cumulatifs expliquent ce ph\'enom\`ene.

Le premier facteur est le \textbf{mismatch linguistique}. Le corpus Wikipedia est en fran\c{c}ais tandis que les questions GeoSQA sont en anglais (traduites du chinois) et les questions GKMC sont en anglais. Ce d\'ecalage emp\^eche le retriever de trouver des documents pertinents (HR@5 = 0.060 \`a 0.075). La validation avec le corpus EN (HR@5 $\times$2.5 \`a $\times$2.9) confirme que cette limitation est surmontable en alignant la langue du corpus sur celle des questions.

Le deuxi\`eme facteur est le \textbf{contexte non pertinent}. Quand le retriever retourne des documents hors sujet, le contexte inject\'e dans le prompt d\'egrade les performances en d\'etournant le mod\`ele de ses connaissances internes. Ce ph\'enom\`ene est coh\'erent avec les observations de Gao et al.~\cite{gao2023ragsurvey} sur la sensibilit\'e du RAG na\"if \`a la qualit\'e du retrieval.

Le troisi\`eme facteur est l'\textbf{exception Tourism-QA}. Sur ce dataset, le corpus est en anglais et sp\'ecifique aux POI \'evalu\'es, le retrieval atteint un Hit Rate de 0.48 (contre 0.060-0.075 sur GeoSQA/GKMC). M\^eme dans cette configuration favorable, seul Qwen3-0.6B b\'en\'eficie marginalement du contexte (+1.8\,\%), tandis que les mod\`eles Llama sont fortement d\'egrad\'es (-21 \`a -23\,\%), confirmant que le RAG na\"if est insuffisant pour am\'eliorer des mod\`eles poss\'edant d\'ej\`a des connaissances internes sur la t\^ache.


\subsection{Pourquoi Qwen3-0.6B domine-t-il sur Tourism-QA en mode RAG ?}

Ce résultat contre-intuitif - le plus petit modèle surpasse les plus grands en 
mode RAG - s'explique par un compromis fondamental du RAG naïf entre connaissances 
internes et dépendance au contexte. Le tableau~\ref{tab:compromis_tourism} illustre 
ce mécanisme.

\begin{table}[H]
\centering
\small
\caption{Illustration du compromis RAG sur Tourism-QA : une baseline élevée 
implique que le RAG dégrade ; une baseline faible implique que le RAG améliore.}
\label{tab:compromis_tourism}
\begin{tabular}{lrrrp{5.5cm}}
\toprule
\textbf{Modèle} & \textbf{Base. Acc@1} & \textbf{RAG Acc@1} & 
\textbf{$\Delta_{rel}$} & \textbf{Interprétation} \\
\midrule
Llama-3.2-3B & 0.710 & 0.590 & $-16.9$\,\% & 
    Bonnes connais. internes $\Rightarrow$ RAG perturbe \\
Qwen3-1.7B   & 0.580 & 0.530 & $-8.6$\,\%  & 
    Connais. modérées $\Rightarrow$ RAG gêne \\
Llama-3.2-1B & 0.440 & 0.410 & $-6.8$\,\%  & 
    Peu de connais. $\Rightarrow$ RAG aide peu \\
Qwen3-0.6B   & 0.600 & 0.660 & $+10.0$\,\% & 
    Dépend du contexte $\Rightarrow$ RAG aide \\
\bottomrule
\end{tabular}
\end{table}

Les modèles disposant de bonnes connaissances internes sur la tâche 
(Llama-3.2-3B, baseline $= 0.71$) subissent les dégradations les plus fortes : 
le contexte des reviews récupérées entre en conflit avec leurs représentations 
internes et les détourne de la bonne réponse. À l'inverse, Qwen3-0.6B, dont les 
connaissances sur le tourisme sont limitées, s'appuie pleinement sur le contexte 
fourni --- ce qui est précisément l'objectif fondateur du RAG selon ~\cite{lewis2020retrieval}. Ce résultat valide partiellement le principe du RAG tout en 
confirmant sa limite principale : il n'est bénéfique que pour les modèles qui en 
ont réellement besoin, c'est-à-dire ceux dont les connaissances internes sur la 
tâche ciblée sont insuffisantes.

\subsection{Sensibilit\'e au contexte r\'ecup\'er\'e sur Tourism-QA}

Tourism-QA r\'ev\`ele un ph\'enom\`ene remarquable. Les mod\`eles poss\'edant de bonnes connaissances internes sont massivement d\'egrad\'es par le RAG. Llama-3.2-3B atteint 0.70 en Baseline mais chute \`a 0.54 en RAG (-22.9\,\%), et Llama-3.2-1B passe de 0.52 \`a 0.41 (-21.2\,\%). \`A l'inverse, les mod\`eles Qwen r\'esistent au contexte. Qwen3-1.7B maintient ses performances (0.59 dans les deux conditions) et Qwen3-0.6B est l\'eg\`erement am\'elior\'e (+1.8\,\%). Ce comportement sugg\`ere que la famille Qwen filtre mieux le contexte non pertinent que la famille Llama, et valide partiellement le principe fondateur du RAG~\cite{lewis2020retrieval}, \`a savoir l'apport de connaissances externes aux mod\`eles qui en manquent.

\subsection{Limites de l'\'etude}

\begin{itemize}[leftmargin=*]
  \item \textbf{\'Echantillon de mod\`eles limit\'e.} Seuls 4 mod\`eles ont produit des r\'esultats (Ministral incompatible), tous de taille inf\'erieure ou \'egale \`a 3B. Les mod\`eles optionnels de 8B n'ont pas pu \^etre charg\'es sur le GPU T4. Les conclusions ne se g\'en\'eralisent pas aux mod\`eles de 7B et plus.

  \item \textbf{Corpus monolingue.} Le corpus Wikipedia FR est inadapt\'e pour les questions en anglais. La validation avec le corpus EN (HR@5 $\times$2.5 \`a $\times$2.9) confirme que cette limitation est surmontable.

  \item \textbf{RAG na\"if uniquement.} Seule la configuration \textit{retrieve-then-read} a \'et\'e test\'ee, sans \'etapes de re-ranking ou de r\'e\'ecriture de requ\^ete. Des techniques de RAG avanc\'e~\cite{gao2023ragsurvey} pourraient am\'eliorer significativement les performances.

  \item \textbf{Variabilit\'e des r\'esultats.} Les \'evaluations portent sur 200 questions pour GeoSQA/GKMC et 100 pour Tourism-QA (seed = 42), ce qui implique une certaine variance d'\'echantillonnage.

  \item \textbf{M\'etriques h\'et\'erog\`enes.} GeoSQA et GKMC utilisent l'Accuracy MCQ tandis que Tourism-QA utilise l'Accuracy@1. La comparaison directe entre datasets reste approximative.
\end{itemize}

\subsection{Note sur la reproductibilité des résultats}

Les valeurs numériques présentées dans ce rapport final proviennent 
d'exécutions définitives avec \texttt{seed = 42} sur les échantillons de 
référence (200 questions pour GeoSQA/GKMC, 100 pour Tourism-QA). De légères 
différences avec les valeurs rapportées dans les rapports intermédiaires de 
Phase 1 et Phase 2 s'expliquent par l'utilisation de sous-ensembles de questions 
différents lors des exécutions exploratoires préliminaires. Le tableau~\ref{tab:repro} 
récapitule les écarts les plus visibles.

\begin{table}[H]
\centering
\caption{Écarts entre les valeurs des rapports intermédiaires (Phase 2) et 
le rapport final, dus à des sous-ensembles d'évaluation différents.}
\label{tab:repro}
\begin{tabular}{llcc}
\toprule
\textbf{Modèle} & \textbf{Condition} & 
\textbf{Valeur Phase 2} & \textbf{Valeur finale} \\
\midrule
Qwen3-0.6B   & Tourism-QA RAG Acc@1  & 0.6600 & 0.5700 \\
Qwen3-1.7B   & Tourism-QA Base Acc@1 & 0.5800 & 0.5900 \\
Llama-3.2-3B & GKMC Base Acc         & 0.5200 & 0.5650 \\
\bottomrule
\end{tabular}
\end{table}

Ces écarts, bien que non négligeables, ne remettent pas en cause les tendances générales observées : l'avantage de la famille Qwen sur Llama, l'impact négatif moyen du RAG naïf, et la supériorité du corpus anglophone pour le retrieval restent stables quelle que soit l'exécution considérée.

% ====================================================================
% 9. CONCLUSION ET PERSPECTIVES
% ====================================================================
\section{Conclusion et perspectives}
\label{sec:conclusion}

\subsection{Synth\`ese des r\'esultats}

Ce projet a permis d'impl\'ementer et d'\'evaluer un syst\`eme de question-r\'eponse g\'eographique en configuration RAG, conform\'ement au cahier des charges. Le syst\`eme a \'et\'e \'evalu\'e avec 4 LLMs de deux familles (Qwen, Llama), sur 3 benchmarks g\'eographiques (GeoSQA, GKMC, Tourism-QA), en comparant syst\'ematiquement les configurations Baseline (sans RAG) et RAG (avec contexte r\'ecup\'er\'e).

\begin{table}[H]
\centering
\caption{Bilan synth\'etique des principaux r\'esultats mesur\'es.}
\label{tab:bilan_final}
\begin{tabularx}{\textwidth}{@{} l X @{}}
\toprule
\textbf{Constat} & \textbf{Valeur mesur\'ee} \\
\midrule
Meilleur mod\`ele global & Qwen3-1.7B (RAG Acc moy. = 0.5100 sur 3 datasets) \\
\rowcolor{lightgray}
Meilleure efficacit\'e (Acc/Taille) & Qwen3-0.6B (0.692 Acc/B) \\
Impact RAG moyen & $\bar{\Delta}_{rel} = -9.4\,\%$ (n\'egatif en moyenne) \\
\rowcolor{lightgray}
Meilleure config. RAG & Qwen3-1.7B + GKMC (+3.4\,\%) \\
Pire d\'egradation RAG & Llama-3.2-1B + GeoSQA (-25.5\,\%) \\
\rowcolor{lightgray}
Corr\'elation taille/performance & $r = 0.19$, $p = 0.56$ (non significatif) \\
Taille aide (intra-famille) & 6/6 configurations \\
\rowcolor{lightgray}
Avantage famille Qwen vs Llama & +0.095 en moyenne (constant sur les 3 datasets) \\
HR@5 (GeoSQA/GKMC, corpus FR) & 0.060 \`a 0.075 (mismatch FR/EN) \\
\rowcolor{lightgray}
HR@5 (Tourism-QA, corpus EN) & 0.48 (corpus EN align\'e) \\
Dataset le plus difficile & GeoSQA (RAG Acc moy. = 0.2575) \\
\rowcolor{lightgray}
Dataset le plus facile & Tourism-QA (RAG Acc moy. = 0.5275) \\
Ministral-3B & Incompatible \texttt{AutoModelForCausalLM} \\
\bottomrule
\end{tabularx}
\end{table}

Les principaux enseignements sont les suivants\,:

\begin{enumerate}[leftmargin=*]
  \item \textbf{Le RAG na\"if d\'egrade les performances en moyenne.} Sur les 12 configurations test\'ees, seules 3 b\'en\'eficient marginalement du RAG ($\bar{\Delta}_{rel}$ = -9.4\,\%). La condition n\'ecessaire est un alignement linguistique et th\'ematique entre le corpus et les questions.

  \item \textbf{La qualit\'e du retrieval est le facteur d\'eterminant.} Le passage d'un Hit Rate de 0.060 (corpus FR) \`a 0.48 (corpus EN sur Tourism-QA) r\'eduit la d\'egradation induite par le RAG, mais ne suffit pas \`a garantir une am\'elioration syst\'ematique.

  \item \textbf{La famille du mod\`ele importe davantage que la taille seule.} L'avantage de Qwen sur Llama est constant (+0.095 en moyenne) et ind\'ependant du dataset. La taille aide dans les 6 comparaisons intra-famille.

  \item \textbf{La famille Qwen r\'esiste mieux au contexte non pertinent que Llama}, comme en t\'emoigne Tourism-QA o\`u les mod\`eles Llama perdent plus de 21\,\% tandis que Qwen reste stable.
\end{enumerate}

\subsection{Perspectives}

Le tableau~\ref{tab:facteurs_phase3} synthétise les facteurs identifiés comme 
déterminants et les actions prioritaires recommandées pour la Phase~3.

\begin{table}[H]
\centering
\small 
\caption{Synthèse des facteurs déterminants identifiés en Phase~2 et actions 
recommandées pour la Phase~3.}
\label{tab:facteurs_phase3}
\renewcommand{\arraystretch}{1.3}
\begin{tabular}{p{3.8cm}lp{3.5cm}p{4cm}}
\toprule
\textbf{Facteur} & \textbf{Impact} & \textbf{Valeur mesurée} & 
\textbf{Action Phase 3} \\
\midrule
Mismatch linguistique FR/EN & Critique & 
    HR@5 : $0.065 \to 0.175$ ($\times 2.9$) & 
    Adopter un corpus exclusivement anglophone \\
Alignement corpus-tâche & Critique & 
    HR@5 : $0.075 \to 0.480$ ($\times 6.4$) sur Tourism-QA & 
    Corpus thématiquement ciblé par dataset \\
Famille (Qwen vs Llama) & Important & 
    Avantage Qwen constant : $+0.087$ en moyenne & 
    Tester Mistral-7B en quantification 4-bit \\
Taille intra-famille & Important & 
    Aide dans 5/6 configurations & 
    Évaluer modèles 7B+ si ressources disponibles \\
RAG naïf vs avancé & Modéré & 
    $\bar{\Delta}_{rel} = -5.5$\,\% & 
    Implémenter re-ranking cross-encoder \\
Connaissances internes & Modéré & 
    RAG utile uniquement pour petits modèles sur tâches ciblées & 
    Cibler le RAG sur modèles $< 1$B \\
\bottomrule
\end{tabular}
\end{table}

\medskip

Plusieurs pistes d'am\'elioration se d\'egagent de cette \'etude.

\begin{itemize}[leftmargin=*]
  \item \textbf{Corpus multilingue.} Remplacer le corpus Wikipedia FR par un corpus en anglais pour \'eliminer le mismatch linguistique qui limite le Hit Rate@5 \`a 0.060-0.075 sur GeoSQA et GKMC.

  \item \textbf{Am\'elioration du retriever.} Tester des strat\'egies de reranking (cross-encoder) et de chunking adaptatif pour am\'eliorer la pertinence au-del\`a du Hit Rate de 0.48 obtenu sur Tourism-QA.

  \item \textbf{Troisi\`eme famille de mod\`eles.} Tester \texttt{Mistral-7B-Instruct-v0.3} en quantification 4-bit via \texttt{bitsandbytes} afin de valider si l'avantage de Qwen sur Llama est g\'en\'eralisable \`a d'autres familles.

  \item \textbf{Mod\`eles de plus grande taille.} \'Evaluer des mod\`eles de 7B \`a 8B param\`etres (Qwen3-8B, Llama-3.1-8B-Instruct) pour v\'erifier si les conclusions sur les petits mod\`eles se transf\`erent.

  \item \textbf{RAG avanc\'e.} Impl\'ementer des techniques de RAG avanc\'e~\cite{gao2023ragsurvey} telles que la r\'e\'ecriture de requ\^ete, le re-ranking par cross-encoder, ou le RAG it\'eratif.
\end{itemize}


% ====================================================================
% REFERENCES
% ====================================================================
\newpage
\bibliographystyle{plain}
\bibliography{references}

\end{document}